
\section{Discussion}
\label{sec:discussion}

\subsection{What the Experiment Establishes}
\label{sec:disc:what}

Three things are established on this benchmark, and they should be stated with
their scope attached. First, the portfolio is complementary: three of four
policies are strictly best somewhere with a margin that survives replication, and
73.0\,\% of contexts have more than one non-dominated policy on the criterion
vector. Second, the decision value of that complementarity is 0.45\,s per vehicle
on a 474.4\,s mean, against 220.2\,s for the first-order choice of which fixed
policy to adopt. Third, a selector trained on decision cost recovers 55.8\,\% of
the headroom on held-out contexts and reduces the worst case by more than half,
while remaining worth 0.25\,s per vehicle.

What is \emph{not} established is anything about a real network. One synthetic
three-corridor topology, one origin--destination structure, one signal plan
family and one demand profile were studied. The numbers above are properties of
that benchmark.

\subsection{Why Complementarity Does Not Imply Decision Value}
\label{sec:disc:gap}

The gap between the second and third paragraphs above is the paper's main point,
and it has a mechanism rather than being an accident of this environment.

Complementarity is a statement about the \emph{argmin} of a cost function over a
portfolio; decision value is a statement about the \emph{gap} between the best
fixed member and that argmin. The two coincide only when the best fixed policy is
badly wrong where it is wrong. Here it is not: P3 attains the minimum in two
thirds of contexts, and in the third where it does not it loses 1.32\,s on average
while winning by 22.47\,s where it does. An upside-to-downside ratio of 17 means
the fixed policy absorbs most of the variation in the winner map, leaving
adaptation to compete for a residue.

This has a practical reading that is easy to miss. The right question to ask of a
portfolio is not ``do different members win in different instances?'' --- in any
non-degenerate portfolio some will --- but ``how much does the best fixed member
lose where it loses, and how often?'' Those two quantities, not the count of
distinct winners, determine whether per-instance selection is worth building.
They are cheap to compute from the same data that establishes complementarity,
and reporting only the first while implying the second overstates the case for
adaptation.

The multi-criteria result reinforces the same point from the other side. On the
criterion vector the portfolio looks strongly complementary --- 73.0\,\% of
contexts admit more than one non-dominated policy, and the reliability-aware
policy is the best stopped-delay policy in two thirds of contexts while almost
never being the best journey-time policy. None of that changes the decision value
on the decision criterion. A Pareto set can be large because criteria disagree
about an ordering while every difference is small.

\subsection{Mechanistic Reading of the Boundary}
\label{sec:disc:mech}

Where the preferred policy does change, the organising quantity is the product of
demand and penetration relative to the capacity of the corridor every policy would
otherwise prefer. Below a controlled-flow saturation of roughly 0.76 the guided
cohort is small enough, or the arterial large enough, that diverting traffic costs
more than the congestion it avoids, and the reliability-aware and reactive
policies hold much of the space. Above it the arterial cannot absorb the guided
cohort and load balancing wins almost everywhere. The quantity is dimensionless
and computable before any vehicle departs, which is what makes it usable as a
switching criterion.

Three qualifications belong with that reading. The threshold is a bracket between
adjacent sampled levels, not a number: a grid cannot locate a boundary more finely
than its own spacing, and the refinement campaign showed that at 150\,veh/h
spacing a third of the transitions are not even monotone in demand. The rule is an
association within a designed factorial; the experiment contains no treatment that
manipulates $x_8$ while holding the physical routes fixed, so it does not identify
a mechanism. And the same rule, evaluated as a decision rule rather than as a
description, is worse than doing nothing --- which is the subject of the next
subsection.

\subsection{Why Decision-Oriented Evaluation Changes the Answer}
\label{sec:disc:decision}

The mechanistic rule is more accurate than the fixed policy and costs more than
it. The most accurate learned rule is not the cheapest. The cheapest rule is the
only one trained on advantage rather than on the winner's identity. Taken
together these say that in this problem accuracy and decision quality are not
merely imperfectly correlated; they order the candidate rules differently.

The reason is structural rather than incidental. The distribution of regret is
extremely skewed: in most contexts every policy is within a fraction of a second
of every other, and in a small minority the gap is tens of seconds. A rule fitted
to the winner's label spends its capacity on the many contexts where the label is
almost arbitrary, and pays for its errors in the few where it is not. A rule fitted
to advantage sees the skew directly: it is indifferent where the advantage is
small, which is most places, and deviates only where it is large. That is why B4
alone improves the tail --- worst case from 17.19\,s to 6.96\,s, CVaR$_{90}$ from
3.67\,s to 1.66\,s --- and why it does so while being \emph{less} accurate than the
rule it beats.

For practice this means that selection accuracy is the wrong acceptance criterion
for a routing-policy selector, and that a selector should be trained on the cost
of its decisions when those costs are available, which in a simulation study they
always are. For this literature it means that reported accuracies of
instance-selection systems are not comparable across problems whose regret
distributions differ in skew.

\subsection{Distribution Shift as a Deployment Problem}
\label{sec:disc:shift}

The shift results are not a benchmark of a learning method; they are a statement
about what an operator would face. Three points follow from them.

Shift is not one thing. The same selector helps on an interpolation control and a
domain shift, is indistinguishable on three factor-level extrapolations, and
harms on two extrapolations and a concept shift. Reporting a single
out-of-distribution number would have described none of these cases, and would
have hidden the one that matters most: the concept shift, where the relationship
between context and policy performance itself changes, is where the damage is
worst, at more than three times the fixed policy's regret.

Support-based detection is bounded by its representation. The gate flagged every
held-out context of the three factor-level extrapolations --- exactly what it was
built to do --- and none of the domain shift, because the descriptor set encodes
whether a disruption may occur but not where. This is not a defect of the
particular gate. Any detector that measures distance in a feature space is blind
to a change the feature space does not encode, and the changes an operator most
needs to catch are often of that kind, because they are the ones nobody
anticipated when the features were designed. It follows that an out-of-
distribution guard should be specified together with an explicit list of the
changes it can and cannot see, and that this list is a property of the features
and not of the detector.

The failed bypass-disruption campaign is relevant here rather than merely
procedural. The strongest available domain-shift test could not be executed
because closing a lane on a single-lane corridor disconnects it instead of
constraining it. The replacement test is real but weaker, and the consequence is
that the study's evidence on domain shift rests on the corridor that no policy
specialises in.

\subsection{Selective Prediction and the Effect-to-Noise Ratio}
\label{sec:disc:selective}

The selective rule never acted. That result is worth more than a working gate
would have been, because it identifies a precondition that is rarely checked.

A confidence-gated selector certifies a deviation when a lower bound on the
predicted advantage is positive. The width of that bound is set by the residual
dispersion of the prediction problem, which here is governed by the replication
noise of the simulation: a mean $2\,$SE noise floor of 20.90\,s against a mean
advantage of 0.45\,s. No calibration of the interval can make it narrower than the
noise it is estimated from, so the gate cannot certify anything, and the
in-distribution gain is forgone in full.

The design rule is therefore to compute the ratio of the available effect to the
replication noise \emph{before} building an abstention mechanism. Where that ratio
is of the order seen here, a confidence gate is not conservative; it is inert, and
the honest options are to accept the unguarded selector with its shift behaviour
documented, to raise the effect size by changing the portfolio, or to reduce the
noise by replication --- which here would require about 25 seeds per context, five
times the campaign that was run.

\subsection{Implications for Routing-System Design}
\label{sec:disc:design}

The results support a small number of concrete statements, all scoped to this
environment.

The first-order decision is which fixed policy to adopt, and it is worth roughly
490 times the second-order decision of whether to adapt it. An operator with a
limited engineering budget should spend it on that choice, on the calibration of
the chosen policy's parameters --- where loosening one tolerance was worth 37
times the entire adaptation headroom --- and on measurement, rather than on an
adaptive layer.

If an adaptive layer is nonetheless built, three things follow from the results.
It should be trained on decision cost rather than on the winner's identity,
because that is what separates the rule that helps from the rule that hurts. It
should be given an explicitly declared operating domain, because the harm under
shift is concentrated in extrapolation above the training range and in concept
change. And its guard should not be assumed to detect shifts that its features do
not encode.

Finally, the reliability-aware policy is worth retaining even though it almost
never minimises journey time, because it minimises stopped delay in two thirds of
contexts. Which of those an operator wants is a policy question and not a
technical one, but the choice should be made explicitly rather than inherited from
whichever criterion the evaluation happened to scalarise.

\subsection{Limitations and Threats to Validity}
\label{sec:limits}

\paragraph{External validity} One synthetic network with three parallel routes,
one origin--destination structure, a homogeneous fleet, and specified rather than
field-calibrated demand and signal plans. Corridor capacities were measured within
the simulation, not from field data. Nothing here establishes the numerical values
for a real network, and the three-corridor structure in particular makes the
candidate path set complete, which is a simplification that removes the
path-generation problem entirely. \emph{This does not invalidate} the separation
between complementarity, decision value and deployability, which is a relationship
between quantities rather than a value of any of them, but it does mean the
magnitudes are benchmark-specific.

\paragraph{Behavioural assumptions} Guided vehicles comply fully with the policy;
unguided vehicles follow a fixed habitual choice; routing is pre-trip with no
en-route re-routing; there is no day-to-day learning and no response to the
guidance system's own past behaviour. Real compliance is partial and
adaptive \citep{toso2024detrimental}, and a system that changes its policy would
face travellers who respond to that change. \emph{This does not invalidate} the
within-benchmark comparisons, which are paired and hold behaviour fixed across
policies, but it does mean the penetration results describe an idealised cohort.

\paragraph{Portfolio scope} Four policies, with parameters frozen before the
study. The sensitivity experiment shows the frozen tolerance of P3 to be
conservative by 16.66\,s per vehicle, so the benchmark understates what that
policy family can do, and a differently tuned portfolio would have a different
headroom. Eco-routing was excluded on screening evidence specific to this network.
\emph{This does not invalidate} the reported headroom, which is a property of the
portfolio as specified, but it does mean the headroom is not an upper bound over
all portfolios.

\paragraph{Statistical resolution} Five paired seeds per context. Under the
declared criterion 60.3\,\% of winning margins are unresolved, and only 24.8\,\% of
the contexts carrying headroom carry more than their own noise band. Resolving the
campaign-wide mean would need about 25 seeds. \emph{This does not invalidate} the
resolved findings --- and the resolved margins are large, with a median of 5.2\,\%
--- but it does mean that fine structure in the winner map is beyond this
experiment's resolution, which is why boundaries are reported as brackets.

\paragraph{Simulation and measurement} Emissions come from a model, not
measurement, and the three criteria are computed from one simulator's outputs.
Journey time includes insertion delay, which is the right convention for a policy
that can oversaturate the network but makes absolute values at the highest demand
partly a statement about the insertion queue. \emph{This does not invalidate} the
policy ordering, which is unaffected by the accounting boundary, but absolute
values should not be compared across studies.

\paragraph{Protocol deviations} Two are recorded in Section~\ref{sec:integrity}:
the main campaign was run after a pre-declared stop gate failed, and the
domain-shift split was replaced after the original proved infeasible. The first is
mitigated by the confirmatory gate declared in advance and by the constraints
adopted at the same time, all of which the results honour; it nonetheless means
the study proceeded past a rule it had set itself. The second means the
domain-shift evidence is weaker than designed. \emph{Neither invalidates} a
reported result, and both are stated so that a reader can discount them as they
see fit.

\paragraph{Causal identification} The experiment is a factorial over context
variables, not an intervention on mechanisms. Every mechanistic statement in
Section~\ref{sec:disc:mech} is an association consistent with the data and is
labelled as such.

\subsection{Research Implications}
\label{sec:disc:research}

Three lines follow from this study rather than being asserted by it. The first is
whether the complementarity--value gap survives in networks whose path sets are
not complete and whose alternatives are less symmetric; the mechanism identified
here --- a near-dominant fixed policy with a large upside-to-downside ratio --- is
not obviously special to three corridors, but nothing here shows that. The second
is whether decision-cost training retains its advantage when the regret
distribution is less skewed, since the argument for it given in
Section~\ref{sec:disc:decision} depends on that skew. The third is the design of
out-of-distribution guards whose coverage is stated in terms of what the feature
representation encodes, which the O8 result suggests is the binding constraint
rather than the calibration method.
